\section{Documentación}


Un apartado que no muchas veces se comenta es la forma de documentación.

La forma más utilizada para poder documentar VHDL es utilzando Doxygen.

\begin{quote}
    Doxygen es una herramienta Open Source utilizada en la documentación de 
    un motón de lenguajes como C, C++, VHDL, etc.
    Para hacer esta documentación se utilizan los comentarios de código.
\end{quote}

Para poder documentar con Doxygen es necesario utilizar el formato en los comentarios
precedidos por '--!'.

\begin{verbatim}
--! Esto es un comentario en formato Doxygen
\end{verbatim}
